% ---------------------------------------------------------------------
%  Chapter 7 : The Propagator
% ---------------------------------------------------------------------
\chapter{The Propagator}\label{ch:propagator}

Every Feynman diagram in this framework is a graph whose \emph{edges} are
copies of one object: the free two-point function, the \term{propagator}
$\Gret$. Wire vertices together with copies of $\Gret$, integrate over the
internal times, and you have a diagram. Nothing downstream --- not the
enumeration, not the type assignment, not a single loop integral --- can run
until $\Gret$ exists. This chapter is about how the engine builds it.

The story has a clean shape. The Gaussian (free) part of the \msrjd{} action is
a bilinear form coupling each response field to each physical field; collect its
coefficients into a matrix $K$ (the ``kinetic'' or ``kernel'' matrix), and the
propagator is simply its inverse $G = K^{-1}$. Because $K$ contains time
derivatives, the natural place to invert it is \emph{frequency space}: Fourier
transform in time to get $K_{\mathrm{ft}}(\omega)$, invert the small matrix once,
and read off a rational function of $\omega$. The physics is then entirely in
that rational function's \emph{poles} (which set relaxation rates), its
\emph{residues} (which set amplitudes), and one \emph{instantaneous} piece that
becomes a $\delta(t)$ in the time domain. The closed-form inverse Fourier
transform reassembles those three ingredients into the time-domain retarded
propagator that the diagram integrators actually call.

This is phases \texttt{[2/7]} (\emph{Build propagator}) and \texttt{[4/7]}
(\emph{Numerical poles + residues}) of the seven-phase trace you watched in
Chapter~\ref{ch:quickstart}. We take them together because they are two halves
of one object: phase~2 builds the \emph{symbolic} frequency-domain propagator
(parameter-independent, cached to disk); phase~4 substitutes the numeric
parameters and extracts the poles and residues. We will build the mathematics
from the ground up, then walk the two source files that implement it ---
\file{pipeline/\_propagator.py} (the engine) and
\file{msrjd/integration/time\_domain/propagator\_td.py} (the time-domain
assembler) --- quoting real code with file and line numbers throughout.

\begin{gotcha}
There is a file named \file{msrjd/core/propagator.py}. It is a \emph{stub} --- a
ten-line docstring saying the logic ``currently lives in the demo notebook.''
Nothing in the working pipeline imports it. The real propagator code is
\file{pipeline/\_propagator.py}. Do not go looking for \code{build\_propagator}
in \file{msrjd/core/propagator.py}; it is not there.
\end{gotcha}

\section{Where the propagator sits}

It helps to fix the data flow before the mathematics. Two functions in
\file{pipeline/\_propagator.py} do the work, and they run at two different points
in \file{pipeline/compute.py}:

\begin{description}
  \item[\code{build\_propagator(ft, model, \ldots)}] (\file{pipeline/\_propagator.py:524})
        Called at phase~2 (\file{pipeline/compute.py:357}). Takes the
        \emph{expanded} \code{FieldTheory} \code{ft} and the \code{model} dict;
        builds the symbolic chain $K_{\mathrm{ker}} \to K_{\mathrm{ft}} \to
        G_{\mathrm{ft}}$, plus the adjugate, the determinant, and the
        instantaneous matrix \code{D\_delta}. It returns a \code{prop} dict with
        the poles and residues still set to \code{None} --- those are
        parameter-dependent and come later. The symbolic stages \emph{are}
        parameter-independent, so this dict is cached to disk.
  \item[\code{compute\_poles\_and\_residues(prop, num\_params, \ldots)}] (\file{pipeline/\_propagator.py:851})
        Called at phase~4 (\file{pipeline/compute.py:632}). Takes the \code{prop}
        dict and \code{num\_params} --- the numeric parameter substitution
        produced by the mean-field solve (phase~3). It substitutes those numbers,
        finds the retarded poles, computes the residue matrices, and
        \emph{mutates} \code{prop} in place to fill \code{prop['pole\_vals']} and
        \code{prop['C\_mats']}.
\end{description}

A third function lives in the other file. \code{build\_G\_t\_matrix}
(\file{msrjd/integration/time\_domain/propagator\_td.py:135}) turns the
frequency-domain data (poles, residues, \code{D\_delta}) into a \emph{time-domain}
matrix $G(t)$, and the per-edge lookups \code{G\_t\_entry} /
\code{G\_t\_delta\_coeff} hand single propagator entries to the diagram
integrators. Schematically:

\begin{lstlisting}[style=console]
  FieldTheory (expanded action)        model dict (params, kernel images)
           |                                       |
           +------------------+--------------------+
                              v
  build_propagator():   K_ker -> K_ft -> G_ft / adj_ft / D_omega / D_delta
                              |   (cached to saved_theories/<tag>/propagator.sobj)
                              v
  compute_poles_and_residues(prop, num_params):
                              fills prop['pole_vals'], prop['C_mats'], prop['D_delta']
                              v
  build_G_t_matrix(prop, t):  {'smooth': sum_k C_k exp(i p_k t),  'delta': D_delta,  't_var'}
                              v
  per-edge   G_t_entry()  /  G_t_delta_coeff()
                              v
  diagram tree integrators (final_integral.py, grouped_integral.py)
\end{lstlisting}

\begin{note}
There is also a precompute entry point at \file{pipeline/\_precompute.py:168}
that calls \code{build\_propagator} at Taylor order~2 purely to warm the cache.
This works because the propagator depends \emph{only} on the bilinear sector of
the action, which is fully captured at Taylor order $\ge 2$ --- so a single
order-2 precompute fills the cache once and every higher-order run reuses it.
The propagator is, in this precise sense, Taylor-order-independent.
\end{note}

\section{From the action to the kernel matrix \texorpdfstring{$K$}{K}}

A stochastic field theory written in \msrjd{} form has an action $S[\varphi,
\tilde\varphi]$. The fields come in two halves of equal count $n_f$:

\begin{itemize}
  \item \term{physical fields} $\varphi_j$, $j = 1,\dots,n_f$ (a density, a
        membrane potential, \ldots),
  \item \term{response fields} $\tilde\varphi_i$, $i = 1,\dots,n_f$, also called
        \emph{auxiliary} or \emph{tilde} fields --- the conjugate variables that
        the \msrjd{} construction introduces, one per physical field.
\end{itemize}

The \term{free action} is the part of $S$ that is bilinear with exactly one
response leg and one physical leg. In the engine's bookkeeping this is
\emph{bigrade $(1,1)$} --- one tilde factor, one physical factor --- and it is
exactly what \code{ft.free\_action()} returns. Schematically,
\begin{equation}
  S_{\mathrm{free}} \;=\; \sum_{i,j}\int\!\dd t\;
  \tilde\varphi_i(t)\,\hat K_{ij}(\Dt)\,\varphi_j(t),
  \label{eq:freeaction}
\end{equation}
where each $\hat K_{ij}$ is a differential/kernel operator in time. For a plain
Ornstein--Uhlenbeck field $\hat K = \Dt + \mu$, so $S_{\mathrm{free}} \supset
\int \tilde n\,(\Dt n + \mu n)$.

\begin{defn}
The \term{kernel matrix} $K$ is the matrix of bilinear coefficients of
\eqref{eq:freeaction}, with layout \emph{[response row, physical column]}:
\[
  K[i, j] \;=\; \text{coefficient of }(\tilde\varphi_i\,\varphi_j)
  \text{ in } S_{\mathrm{free}}.
\]
The propagator is its inverse, $G = K^{-1}$.
\end{defn}

The code extracts $K$ by reading the coefficient of each $(\tilde\varphi_i\,
\varphi_j)$ monomial out of the free action. \code{build\_propagator} first asks
the \code{FieldTheory} for the free action and the polynomial-ring generator
names, then splits the names into the response half and the physical half
(\file{pipeline/\_propagator.py:586}):

\begin{lstlisting}
S_free = ft.free_action()
ring_gen_names = [str(g) for g in R.gens()]

resp_names = ring_gen_names[:ft._n_tilde]
phys_names = ring_gen_names[ft._n_tilde:]
pos_to_row = {ring_gen_names.index(nm): i for i, nm in enumerate(resp_names)}
pos_to_col = {ring_gen_names.index(nm): j for j, nm in enumerate(phys_names)}
\end{lstlisting}

The first \code{ft.\_n\_tilde} generators are the response fields and the rest are
the physical fields; \code{pos\_to\_row} and \code{pos\_to\_col} map a generator's
position in the ring to its row (if it is a response field) or column (if it is
physical). The extraction loop then walks the \emph{monomials} of the free action
--- \code{S\_free.dict()} returns a mapping from exponent vector to coefficient ---
and for each monomial finds which generators carry a positive exponent
(\file{pipeline/\_propagator.py:595}):

\begin{lstlisting}
nf = len(resp_names)
K_data = [[SR(0)] * nf for _ in range(nf)]
for exp_vec, coeff in S_free.dict().items():
    row = col = None
    for idx in range(len(ring_gen_names)):
        if exp_vec[idx] > 0:
            if idx in pos_to_row:
                row = pos_to_row[idx]
            if idx in pos_to_col:
                col = pos_to_col[idx]
    if row is not None and col is not None:
        K_data[row][col] += SR(coeff)
K_mat = matrix(SR, K_data)
\end{lstlisting}

A bilinear $(1,1)$ monomial has exactly one response generator and one physical
generator with positive exponent, so it pins down exactly one $(row, col)$, and
its coefficient is accumulated into $K[\,row, col\,]$. The result \code{K\_mat} is
a Sage symbolic matrix of size $n_f \times n_f$.

\begin{note}
\code{SR} is SageMath's \term{Symbolic Ring} --- the general symbolic-expression
type, a Python wrapper around the GiNaC/Pynac C++ engine. \code{SR(x)} coerces
anything (an integer, a coefficient, a string) to a symbolic expression, and
\code{matrix(SR, data)} builds a matrix whose entries live in that ring. Almost
every matrix in this chapter lives over \code{SR}. We meet several more Sage
rings below; each is introduced where it first appears.
\end{note}

\section{Kernel form: the \texorpdfstring{$\delta$}{delta} and \texorpdfstring{$\delta'$}{delta-prime} symbols}

We want to Fourier transform $K$ in time, but a raw coefficient like $\Dt n$ (a
time derivative) does not Fourier-transform cleanly as written. The engine first
rewrites every entry of $K$ into \term{kernel form}, a sum of a Dirac delta and
its derivative:
\begin{equation}
  c \;\longmapsto\; c_0\,\delta(t) \;+\; c_1\,\delta'(t),
  \qquad c_0 = c\big|_{\Dt \to 0}.
  \label{eq:kernelform}
\end{equation}
There are three abstract namespace symbols involved, all carried on the
\code{FieldTheory}'s namespace object \code{ns}:

\begin{itemize}
  \item \code{ns.Dt} --- a marker for the time-derivative operator $\Dt$,
  \item \code{ns.delta\_D} --- the Dirac delta $\delta(t)$ (displayed as $\delta$),
  \item \code{ns.delta\_Dp} --- its derivative $\delta'(t)$ (displayed as $\delta'$).
\end{itemize}

The helper \code{\_to\_kernel} (\file{pipeline/\_propagator.py:507}) performs the
rewrite \eqref{eq:kernelform}:

\begin{lstlisting}
def _to_kernel(c, Dt, delta_D, delta_Dp):
    c = SR(c)
    if c.has(delta_D):
        return c
    p0   = c.subs({Dt: 0})
    rest = (c - p0).subs({Dt: delta_Dp})
    return p0 * delta_D + rest
\end{lstlisting}

It splits the coefficient into a constant part \code{p0} (set $\Dt \to 0$) and a
derivative part \code{rest} (the remainder, with $\Dt$ replaced by the symbol
$\delta'$), and returns \code{p0*delta\_D + rest}. So a plain constant $\mu$
becomes $\mu\,\delta(t)$ --- an instantaneous coupling --- and a derivative $\Dt$
becomes $\delta'(t)$.

Why wrap the \emph{constant} part in $\delta(t)$ at all? Because of what the
Fourier transform does to a bare constant. The transform of $\delta(t)$ is the
constant $1$, but the transform of a bare constant $\mu$ is the distribution
$2\pi\,\mu\,\delta(\omega)$ --- not what we want. Writing $\mu$ as $\mu\,\delta(t)$
ensures the transform returns the clean constant $\mu$ back (the docstring at
\file{pipeline/\_propagator.py:514} spells this out). Abstract \emph{kernel}
symbols (for example a synaptic kernel \code{ns.g} in a coloured-noise model) are
not constants and are left untouched here; their frequency images are applied
later, after the transform.

\code{build\_propagator} applies \code{\_to\_kernel} entrywise to build
\code{K\_ker} (\file{pipeline/\_propagator.py:615}):

\begin{lstlisting}
Dt       = ns.Dt
delta_D  = ns.delta_D
delta_Dp = ns.delta_Dp

K_ker = matrix(
    SR,
    [[_to_kernel(K_mat[i, j], Dt, delta_D, delta_Dp)
      for j in range(nf)] for i in range(nf)],
)
\end{lstlisting}

\section{The Fourier convention --- and why poles live in the upper half-plane}

The engine uses a single, fixed Fourier convention pipeline-wide, defined in
\code{fourier\_transform} (\file{msrjd/core/field\_theory.py:33}). It is the
angular-frequency convention with no $2\pi$ in the exponent:
\begin{equation}
  F(\omega) \;=\; \int_{-\infty}^{\infty} f(t)\,\ee^{-\ii\omega t}\,\dd t.
  \label{eq:ftconv}
\end{equation}
Under \eqref{eq:ftconv} the kernel symbols transform as
\[
  \delta(t) \;\longmapsto\; 1, \qquad
  \delta'(t) \;\longmapsto\; \ii\omega,
\]
so a kernel-form entry $c_0\,\delta + c_1\,\delta'$ transforms to $c_0 + c_1\,
(\ii\omega)$. For the OU example $\hat K = \Dt + \mu$ becomes $K_{\mathrm{ft}} =
\ii\omega + \mu$.

\code{build\_propagator} performs the transform by substituting the abstract
$\delta$ symbols for their concrete Sage forms --- \code{dirac\_delta(t)} for
$\delta$ and \code{diff(dirac\_delta(t), t)} for $\delta'$ --- and calling
\code{fourier\_transform} on each nonzero entry
(\file{pipeline/\_propagator.py:624}):

\begin{lstlisting}
t_var = SR.var('t')
omega = SR.var('omega', latex_name=r'\omega')
time_subs = {
    delta_D:  dirac_delta(t_var),
    delta_Dp: diff(dirac_delta(t_var), t_var),
}
K_ft_data = [[SR(0)] * nf for _ in range(nf)]
for i in range(nf):
    for j in range(nf):
        c = K_ker[i, j]
        if not c.is_zero():
            K_ft_data[i][j] = fourier_transform(
                SR(c).subs(time_subs), t_var, omega
            )
K_ft = matrix(SR, K_ft_data)
\end{lstlisting}

After the transform, any abstract kernel symbol still present (the
coloured-noise case) is replaced by its frequency image $\hat g(\omega)$ via a
model-supplied hook (\file{pipeline/\_propagator.py:642}):

\begin{lstlisting}
kft_hook = model.get('kernel_ft_image')
if kft_hook is not None:
    kft_subs = kft_hook(ns, omega)
    K_ft = K_ft.apply_map(lambda e: SR(e).subs(kft_subs))
\end{lstlisting}

The \code{kernel\_ft\_image} hook is a callable the theory file supplies; it maps
abstract kernel symbols to their frequency images (for a white-noise model there
are none, and the hook is absent).

\begin{defn}
\code{dirac\_delta(t)} is Sage's symbolic Dirac delta; \code{diff(expr, var)} is
symbolic differentiation, so \code{diff(dirac\_delta(t), t)} is the symbolic
$\delta'(t)$. \code{fourier\_transform} is the engine's own wrapper
(\file{msrjd/core/field\_theory.py}); internally it dispatches the integral to
SymPy first (which handles \code{positive=True} assumptions on time constants
without stopping to ask the sign of $\omega$) and falls back to Maxima. SymPy is
therefore on the hot path even though no \code{import sympy} appears in the
propagator files.
\end{defn}

\begin{gotcha}
\textbf{Retarded $=$ $\mathrm{Im}(\omega) > 0$ is the single most important sign
convention in this subsystem, and it is born directly from the $\ee^{-\ii\omega
t}$ in \eqref{eq:ftconv}.} Because the forward transform carries
$\ee^{-\ii\omega t}$, the \emph{inverse} transform carries $\ee^{+\ii\omega t}$
and closes its $\omega$-contour in the \emph{upper} half-plane for $t > 0$. So the
retarded (causal) propagator is parameterised by poles with $\mathrm{Im}(\omega) >
0$: each $\ee^{\ii\omega_k t}$ then decays for $t > 0$ and grows for $t < 0$.
Every pole filter in the code keeps roots with \code{imag > 1e-9}. If you ever
flipped \code{fourier\_transform} to $\ee^{+\ii\omega t}$, every pole-side sign
would flip and the entire downstream causality machinery would break. This is
restated in the header of \file{propagator\_td.py:12}.
\end{gotcha}

\section{The frequency-domain propagator \texorpdfstring{$G_{\mathrm{ft}} = K_{\mathrm{ft}}^{-1}$}{G = K-inverse}}

With $K_{\mathrm{ft}}(\omega)$ in hand, the propagator is the matrix inverse
\begin{equation}
  G_{\mathrm{ft}}(\omega) \;=\; K_{\mathrm{ft}}(\omega)^{-1},
  \label{eq:Ginv}
\end{equation}
a matrix of \emph{rational functions of $\omega$}. Two companion objects come for
free and matter later:

\begin{itemize}
  \item the \term{determinant} $D_\omega = \det K_{\mathrm{ft}}$. Its numerator is
        the \emph{characteristic polynomial}, whose roots are the system poles.
  \item the \term{adjugate} $\mathrm{adj}(K_{\mathrm{ft}})$, the transpose of the
        cofactor matrix. By Cramer's rule $K^{-1} = \mathrm{adj}(K)/\det(K)$, so
        $\mathrm{adj}(K_{\mathrm{ft}}) = G_{\mathrm{ft}}\cdot\det(K_{\mathrm{ft}})$.
\end{itemize}

For a diagonally-dominant $K_{\mathrm{ft}}$ every entry of $G_{\mathrm{ft}}$
shares one common denominator $Q(\omega)$, the polynomial numerator of
$\det(K_{\mathrm{ft}})$.

\subsection{The complexity gate: when symbolic inversion is skipped}

Sage's symbolic \code{K\_ft.inverse()} is \emph{cofactor-based}: it expands all
$n_f^2$ cofactor sub-determinants symbolically, an $O(n_f!)$ operation in the
matrix dimension. For a small matrix (the OU model is $1\times1$; a two-field
system is $2\times2$) this is instantaneous. For a large matrix --- a spike-reset
model can reach $n_f = 8$ --- it can run for \emph{hours}, even when the symbol set
is lean, simply because $8! = 40\,320$ minor determinants must each be simplified
as rational symbolic expressions.

So \code{build\_propagator} gates the symbolic inversion on matrix size
(\file{pipeline/\_propagator.py:681}):

\begin{lstlisting}
free_syms = set()
for i in range(nf):
    for j in range(nf):
        free_syms.update(SR(K_ft[i, j]).variables())
free_syms.discard(omega)
rich = nf >= 6 or len(free_syms) > 20
\end{lstlisting}

A matrix is \term{rich} if $n_f \ge 6$ (the dominant trigger) \emph{or} if it
carries more than twenty free symbols beyond $\omega$ (a secondary trigger). The
comment at \file{pipeline/\_propagator.py:672} records the lesson that motivated
the matrix-size metric: spike-reset at $n_f = 8$ with only $\sim18$ free symbols
ran for hours under the old symbol-count gate, because the cost is $O(n_f!)$ in
the \emph{dimension}, not in the symbol count.

\begin{itemize}
  \item If the matrix is \emph{not} rich, the engine attempts the symbolic inverse
        under a wall-clock budget (next subsection).
  \item If it \emph{is} rich, symbolic inversion is skipped entirely:
        \code{G\_ft}, \code{adj\_ft}, \code{D\_omega}, and \code{D\_delta} are all
        left \code{None}, and \emph{everything} is deferred to phase~4, which
        works numerically after the parameters are substituted (at which point
        every entry is a univariate rational in $\omega$, far cheaper to handle).
\end{itemize}

\subsection{The budgeted symbolic-inverse block}

When the matrix is lean, the symbolic inverse, determinant, adjugate, and
\code{D\_delta} are all computed in one block, wrapped in a $120$-second SIGALRM
budget (\file{pipeline/\_propagator.py:706}):

\begin{lstlisting}
def _do_symbolic_inverse_block():
    t0 = _time.perf_counter()
    G    = K_ft.inverse()
    D_om = K_ft.det()
    # Adjugate via the identity  adj(K) = K^(-1) . det(K).
    adj  = G * D_om
    ...
    return G, D_om, adj, matrix(SR, Dd_data)
\end{lstlisting}

Two design choices in this tiny block are worth calling out, because each is a
hard-won performance fix.

\begin{gotcha}
\textbf{The adjugate is computed as \code{G * D\_om}, never via Sage's
\code{.adjugate()}.} Sage's \code{K\_ft.adjugate()} re-expands all $n_f^2$
cofactor sub-determinants from scratch, throwing away the work the inverse
already did. The comment at \file{pipeline/\_propagator.py:710} records the
measurement: $734$~s for \code{.adjugate()} versus $\sim11$~s for the inverse
itself on a $4\times4$ spike-reset matrix. Reusing the inverse via
$\mathrm{adj} = G\cdot\det$ is the whole point.
\end{gotcha}

If the block exceeds its budget --- or raises any exception, or aborts deep inside
a native routine with a \code{SignalError} --- the four symbolic stages are set to
\code{None} and phase~4 picks up the slack numerically
(\file{pipeline/\_propagator.py:731}):

\begin{lstlisting}
try:
    G_ft, D_omega, adj_ft, D_delta = _run_with_timeout(
        _do_symbolic_inverse_block, _SYMBOLIC_INVERSE_BUDGET_SEC)
except _PolynomialPathTimeout:
    ...
    G_ft = adj_ft = D_omega = D_delta = None
\end{lstlisting}

\code{\_run\_with\_timeout} (\file{pipeline/\_propagator.py:57}) is a small helper
that installs a POSIX \code{SIGALRM} handler raising \code{\_PolynomialPathTimeout}
after the budget expires, arms \code{signal.alarm}, runs the function, and cancels
the alarm in a \code{finally}. On a platform without \code{SIGALRM} (Windows) or
when the budget is \code{None}, it simply calls the function directly --- a no-op
timeout. This is the lower-level cousin of the \code{cysignals} alarm we meet
later for \code{factor()}.

The resulting symbolic stages (or their \code{None} placeholders) are assembled
into the \code{prop} dict, together with the as-yet-unfilled poles and residues
(\file{pipeline/\_propagator.py:765}):

\begin{lstlisting}
prop = {
    'K_ker':         K_ker,
    'K_ft':          K_ft,
    'G_ft':          G_ft,
    'adj_ft':        adj_ft,
    'D_omega':       D_omega,
    'D_delta':       D_delta,
    't_var':         t_var,
    'omega':         omega,
    'nf':            nf,
    'ring_gen_names': ring_gen_names,
    # Filled in later by compute_poles_and_residues():
    'pole_vals':     None,
    'C_mats':        None,
}
\end{lstlisting}

\section{Poles as zeros of \texorpdfstring{$\det K_{\mathrm{ft}}$}{det K}}

Phase~4 begins by substituting the numeric parameters. Once that is done, every
entry of $K_{\mathrm{ft}}$ is a univariate rational function in $\omega$ with
numeric coefficients, and the propagator's poles are the roots of the
characteristic polynomial $Q(\omega) = \operatorname{num}(\det K_{\mathrm{ft}})$
that lie in the upper half-plane.

\begin{defn}
A \term{pole} $\omega_k$ is a (retarded) zero of $\det K_{\mathrm{ft}}$ with
$\mathrm{Im}(\omega_k) > 0$. It sets a relaxation or oscillation rate of the free
theory. A \term{simple} pole is a single root; a higher-multiplicity (repeated)
root is a \emph{Jordan-block} pole and is handled differently --- see the
limitations section.
\end{defn}

Let $\omega_k$ be a simple zero of $Q(\omega)$. The residue of $G_{\mathrm{ft}}$
there is, entrywise,
\begin{equation}
  \operatorname*{Res}_{\omega=\omega_k} G_{\mathrm{ft}}(\omega)
  \;=\; \frac{P_{ij}(\omega_k)}{Q'(\omega_k)},
  \qquad
  G_{\mathrm{ft}}[i,j] = \frac{P_{ij}(\omega)}{Q(\omega)}.
  \label{eq:residue}
\end{equation}
The code stores the residue matrix \emph{scaled by $\ii$} --- the ``$C_k$
convention'' --- because that factor makes the inverse Fourier transform come out
clean:
\begin{equation}
  \texttt{C\_mats}[k] \;=\; \ii\cdot
  \operatorname*{Res}_{\omega=\omega_k} G_{\mathrm{ft}}(\omega).
  \label{eq:Ck}
\end{equation}

The residues are \emph{always} computed numerically, even when a symbolic
adjugate is available. This is not laziness; it is precision. The reason gets its
own gotcha below. First, the primary path.

\section{Residues, tier 1: exact arithmetic over \texorpdfstring{$\mathbb{Q}(i)[\omega]$}{Q(i)}}

\code{compute\_poles\_and\_residues} tries three strategies in order. The first,
\code{\_compute\_residues\_via\_polynomial\_fracfield}
(\file{pipeline/\_propagator.py:81}), is the one that runs for almost every
model. It does \emph{exact} arithmetic --- no floating point until the very last
root-find --- over the field $\mathbb{Q}(i)$.

\subsection{Why an exact field}

The trick is the choice of coefficient ring. Sage can invert a matrix over the
field of rational functions $\mathrm{Frac}(R[\omega])$, and when $R$ is an
\emph{exact} field, polynomial GCDs are reliable: Sage returns every entry of the
inverse in genuinely \emph{canonical irreducible} form $P_{ij}/Q$, with no
spurious cancellable factors and no floating-point GCD noise. The denominator is
then \emph{exactly} the characteristic polynomial.

The exact field used is \code{CyclotomicField(4, 'I\_')}, which is
$\mathbb{Q}(\zeta_4) = \mathbb{Q}(i)$ --- the rationals extended by a primitive
fourth root of unity, which is just the imaginary unit $i$.

\begin{defn}
\term{\code{CyclotomicField(4)}} $= \mathbb{Q}(i)$ is the field of Gaussian
rationals $a + b\,i$ with $a, b \in \mathbb{Q}$. \term{\code{PolynomialRing(R, 'x')}}
builds the ring $R[x]$ of univariate polynomials over $R$; calling
\code{.fraction\_field()} on it gives $\mathrm{Frac}(R[x])$, the rational
functions $P/Q$. The point of doing the matrix inversion over an exact field is
that GCD-over-a-field is exact, so each inverse entry is reduced to true
irreducible $P/Q$ form. A floating-coefficient (\code{CDF}$[\omega]$) path, by
contrast, once left ``degree-18 polynomials with $\sim13$ spurious roots''
(\file{pipeline/\_propagator.py:104}).
\end{defn}

\subsection{Walking the exact path}

The function builds the exact rings, substitutes the parameters, converts every
entry into the exact fraction field, and inverts under the same SIGALRM watchdog
(\file{pipeline/\_propagator.py:115}):

\begin{lstlisting}
CF = CyclotomicField(4, 'I_')
i_CF = CF.gen()
PR_exact = PolynomialRing(CF, 'om_exact')
F_exact  = PR_exact.fraction_field()

K_ft_num = K_ft.apply_map(lambda e: SR(e).subs(num_params))
\end{lstlisting}

The two converters are the heart of it. \code{\_sr\_complex\_to\_CF}
(\file{pipeline/\_propagator.py:122}) turns a symbolic complex \emph{constant}
into an exact $\mathbb{Q}(i)$ element by rationalising its real and imaginary
parts, and \code{\_sr\_to\_F\_exact} (\file{pipeline/\_propagator.py:128}) lifts a
whole rational entry into \code{F\_exact} by converting its numerator and
denominator coefficient lists:

\begin{lstlisting}
def _sr_complex_to_CF(c):
    c = SR(c)
    re = QQ(c.real_part())
    im = QQ(c.imag_part())
    return CF(re) + CF(im) * i_CF
\end{lstlisting}

Here \code{QQ} is Sage's field of rational numbers $\mathbb{Q}$; \code{QQ(x)}
rationalises a float. Once every entry is in \code{F\_exact}, the matrix is built
and inverted (\file{pipeline/\_propagator.py:140}):

\begin{lstlisting}
K_ft_F = matrix(F_exact, K_ft_F_data)
G_ft_F = _run_with_timeout(
    K_ft_F.inverse, _POLYNOMIAL_PATH_BUDGET_SEC)
\end{lstlisting}

\begin{gotcha}
The exact path \emph{requires} parameters that \code{QQ} can rationalise. The line
\code{re = QQ(c.real\_part())} throws on a non-rationalisable (for example,
irrational symbolic) parameter, the whole tier returns a triple of \code{None},
and the caller falls through to the numerical cofactor path
(\file{pipeline/\_propagator.py:122}). In practice parameters are floats, which
\code{QQ} rationalises happily, so the exact path is what runs.
\end{gotcha}

\subsection{The universal denominator must be an LCM}

After the inverse, the code needs the characteristic polynomial $Q(\omega)$. It
is tempting to read it off the first nonzero entry's denominator --- but that is
wrong for a structured $K_{\mathrm{ft}}$. The fix is to take the \emph{least
common multiple} of all the per-entry denominators
(\file{pipeline/\_propagator.py:172}):

\begin{lstlisting}
Q_poly = None
for i in range(nf):
    for j in range(nf):
        if G_ft_F[i, j] == 0:
            continue
        d_ij = G_ft_F[i, j].denominator()
        if Q_poly is None:
            Q_poly = d_ij
        else:
            Q_poly = Q_poly.lcm(d_ij)
\end{lstlisting}

\begin{gotcha}
\textbf{The universal denominator is an LCM, not the first entry's denominator.}
An upper-triangular $K_{\mathrm{ft}}$ --- the shape produced by the
Markovian-embedding preprocessor for coloured noise --- gives entries with
\emph{different} canonical denominators: one entry might carry a single pole
factor while another carries the product of two. Picking the first nonzero entry
would under-count the system poles. The LCM is the universal denominator whose
roots include every system pole; a pole that does not divide a given entry's own
denominator simply yields a vanishing residue there
(\file{pipeline/\_propagator.py:157}).
\end{gotcha}

\subsection{The squarefree guard}

Before root-finding, the code checks that $Q(\omega)$ is \emph{squarefree} --- no
repeated roots (\file{pipeline/\_propagator.py:196}):

\begin{lstlisting}
if not Q_poly.is_squarefree():
    if verbose:
        print(...)
    return None, None, None
\end{lstlisting}

A non-squarefree $Q$ signals a higher-multiplicity (Jordan-block) pole, for
example $Q = (\ii\omega + \mu)^2$ from a Markovianised OU process at $\mu =
1/\tau_c$. The single-pole residue formula \eqref{eq:residue} cannot represent the
$\tau\,\ee^{-pt}$ term such a pole produces, so the exact path bails and lets the
fallback tier take over (the fallback shares the limitation but is more robust to
the numerical near-degeneracy). This is the most significant \emph{known
limitation} of the subsystem; the limitations section returns to it.

\subsection{Roots and residues}

With a squarefree $Q$, the polynomial is converted from the exact field to
\code{CDF}$[\omega]$ (floating-point) for numerical root-finding, the retarded roots
are kept and sorted, and the residues are read off entry by entry.

\begin{defn}
\term{\code{CDF}} is Sage's \term{Complex Double Field}: complex numbers backed by
the machine \code{double} type ($\approx 10^{-16}$ precision). It is used for
root-finding (\code{Q\_cdf.roots(CDF)}) and to store the residue matrices as
\code{matrix(CDF, \ldots)}. \code{complex(value)} converts a \code{CDF} number back
to a plain Python \code{complex}.
\end{defn}

The retarded filter and sort (\file{pipeline/\_propagator.py:214}):

\begin{lstlisting}
roots = Q_cdf.roots(CDF)
pole_vals = [complex(r) for r, _ in roots
             if complex(r).imag > 1e-9]
pole_vals.sort(key=lambda p: (p.imag, p.real))
\end{lstlisting}

The residue loop then implements \eqref{eq:Ck} per entry, $C_k[i,j] = \ii\,
P_{ij}(\omega_k)/Q'_{ij}(\omega_k)$, skipping any entry for which $\omega_k$ is
not a pole of that entry's own denominator (\file{pipeline/\_propagator.py:246}):

\begin{lstlisting}
for pk in pole_vals:
    C_entries = [[0j] * nf for _ in range(nf)]
    for i in range(nf):
        for j in range(nf):
            ...
            qij_at = complex(Q_entry_cdf[i][j](pk))
            if abs(qij_at) > 1e-9:
                continue           # pk is not a pole of this entry
            qijp = complex(Q_entry_prime_cdf[i][j](pk))
            if abs(qijp) < 1e-30:
                continue           # higher-multiplicity: leave 0
            P_at = complex(P_cdf[i][j](pk))
            C_entries[i][j] = 1j * P_at / qijp
\end{lstlisting}

The function returns the triple \code{(pole\_vals, C\_mats, D\_delta)}
(\file{pipeline/\_propagator.py:301}); the third element is the instantaneous
matrix, computed in the same pass and discussed shortly. On success,
\code{compute\_poles\_and\_residues} stores the poles and residues and
\emph{returns early}, never touching the fallback tiers
(\file{pipeline/\_propagator.py:892}):

\begin{lstlisting}
pole_vals_tier1, C_mats_tier1, D_delta_tier1 = (
    _compute_residues_via_polynomial_fracfield(
        K_ft, omega, num_params, nf, verbose=verbose))
if pole_vals_tier1 is not None:
    prop['pole_vals'] = pole_vals_tier1
    prop['C_mats']    = C_mats_tier1
    if prop.get('D_delta') is None and D_delta_tier1 is not None:
        prop['D_delta'] = D_delta_tier1
    ...
    return prop
\end{lstlisting}

Note the third clause: if \code{build\_propagator}'s symbolic block left
\code{D\_delta} as \code{None} (because it busted its budget or the matrix was
rich), the exact path's \code{D\_delta} fills it in.

\begin{note}
The function's docstring (\file{pipeline/\_propagator.py:109}) says it returns
\code{(None, None)} on failure, but every actual \code{return} statement in the
function --- success and failure alike --- is a \emph{three}-tuple
\code{(pole\_vals, C\_mats, D\_delta)}. The caller unpacks three values, so the
code is self-consistent; the docstring is simply stale. When code and prose
disagree, the code is authoritative.
\end{note}

\section{Residues, tiers 2 and 3: the numerical cofactor fallback}

When the exact path returns \code{None} --- non-rationalisable parameters, a
non-squarefree denominator, or a timeout --- \code{compute\_poles\_and\_residues}
substitutes the parameters into $K_{\mathrm{ft}}$ and proceeds numerically
(\file{pipeline/\_propagator.py:926} onward). Which numerical branch runs depends
on whether \code{build\_propagator} left a symbolic adjugate behind:

\begin{itemize}
  \item the \emph{rich / heterogeneous} branch fires when \code{adj\_ft} or
        \code{G\_ft} is \code{None} (\file{pipeline/\_propagator.py:939}), and
  \item the \emph{legacy symbolic} branch fires when both are present
        (\file{pipeline/\_propagator.py:1167}).
\end{itemize}

Both branches find poles the same way --- they take the symbolic
$\det(K_{\mathrm{ft,num}})$, extract its numerator's $\omega$-coefficients, build a
characteristic polynomial, and root it. A sanity check guards against a common
silent failure (\file{pipeline/\_propagator.py:971}):

\begin{lstlisting}
free_syms = set()
for i in range(nf):
    for j in range(nf):
        free_syms.update(SR(K_ft_num[i, j]).variables())
free_syms.discard(omega)
if free_syms and verbose:
    print(f'      WARNING: K_ft_num has free symbols beyond omega: '
          f'{sorted(str(s) for s in free_syms)}.  Pole-finding '
          f'will likely return 0 roots.  ...')
\end{lstlisting}

\begin{gotcha}
\textbf{Pole-finding silently returns zero roots if any free symbol other than
$\omega$ survives.} If a kernel symbol was not substituted (a
\code{kernel\_ft\_image} hook missing an entry) or a saddle value is absent from
\code{num\_params}, then $\det(K_{\mathrm{ft,num}})$ has non-numeric coefficients
and the root-finder finds nothing. The rich branch \emph{warns} about this
explicitly (\file{pipeline/\_propagator.py:976}); the empty pole list then
surfaces downstream as a warning inside \code{build\_G\_t\_matrix}. If you ever
see ``$0$ retarded poles'' on a model you expect to have some, an
unsubstituted symbol is the first thing to check.
\end{gotcha}

Both branches cache $K$'s entrywise numerator/denominator coefficients once
(\file{pipeline/\_propagator.py:1035}) so that the kernel matrix can be evaluated
at any $\omega$ in microseconds via Horner's rule (\code{\_horner},
\file{pipeline/\_propagator.py:841}), candidates are deduplicated, and each is
\term{Newton-refined} on the numerical $\det(K(\omega))$ --- Newton's method
polishing each root to machine precision, with acceptance gated on
$|\det(K(\omega))|/\text{scale} < 10^{-9}$, staying in the upper half-plane, and
deduplication.

The residue evaluation is the common, important part. In \emph{both} branches the
residue is computed via the numerical \term{cofactor adjugate}, never the
symbolic \code{adj\_ft} (\file{pipeline/\_propagator.py:1470}):

\begin{lstlisting}
def _cofactor_adj(M):
    """adj(M) = transpose of cofactor matrix.
    adj[i,j] = (-1)^(i+j) . det(M with row j and col i removed)."""
    n = M.shape[0]
    A = np.zeros_like(M)
    for i in range(n):
        for j in range(n):
            minor = np.delete(np.delete(M, j, axis=0), i, axis=1)
            A[i, j] = ((-1) ** (i + j)) * np.linalg.det(minor)
    return A
\end{lstlisting}

At a pole $K(\omega_k)$ is singular, so $K^{-1}$ does not exist --- but the
adjugate does, by direct cofactor expansion. The determinant derivative
$Q'(\omega_k)$ is taken by central difference, and the residue follows
(\file{pipeline/\_propagator.py:1482}):

\begin{lstlisting}
for pk in pole_vals:
    K_at_pk = _K_at(pk)
    adj_at_pk = _cofactor_adj(K_at_pk)
    h = 1e-6 * (1 + abs(pk))
    d_prime = (_det_at(pk + h) - _det_at(pk - h)) / (2 * h)
    ...
    C_np = 1j * adj_at_pk / d_prime
\end{lstlisting}

\begin{gotcha}
\textbf{The cofactor index convention is transposed.} \code{\_cofactor\_adj} sets
\code{adj[i,j] = (-1)\^{}(i+j) * det(M with row j and col i removed)} --- note
row~$j$, column~$i$. The adjugate is the transpose of the cofactor matrix, and
getting this backwards would transpose every residue matrix
(\file{pipeline/\_propagator.py:1470}).
\end{gotcha}

\begin{gotcha}
\textbf{Residues are always numerical, even when a symbolic \code{adj\_ft}
exists.} Sage stores the symbolic adjugate as a sum of rationals whose per-term
denominators diverge by factors of $10^3$ to $10^4$ as $\omega$ approaches a
kernel pole. Evaluating that sum at the pole causes \term{catastrophic
cancellation} --- nearly-equal large terms subtracting to lose $3$ to $5$
significant digits (a measured $11\%$ to $42\%$ relative error on a quartic
model's residues). The numpy cofactor path builds the minor determinants directly
and has no such leak, so both numerical branches use it
(\file{pipeline/\_propagator.py:1503}). The ``lean symbolic'' branch keeps the
symbolic \code{adj\_ft} around for display, but computes residues numerically all
the same.
\end{gotcha}

\begin{note}
The upshot is that the three tiers are \emph{numerically equivalent} in current
code: the legacy symbolic path was patched to use the numpy cofactor evaluation
too, so tiers~2 and~3 differ only in how they \emph{find the poles}, not in how
they evaluate the residues. One cosmetic redundancy is worth knowing: for the
rich branch there are two nearly-identical Newton-refine blocks
(\file{pipeline/\_propagator.py:1297} and \file{:1389}), so a candidate can be
refined twice --- harmless, since the second refine of an already-converged root
is a near no-op.
\end{note}

\section{The instantaneous limit \texorpdfstring{$D_\delta$}{D-delta}}

If $G_{\mathrm{ft}}(\omega)$ does not vanish as $\omega \to \infty$ --- a
\emph{non-proper} rational function --- then its inverse Fourier transform
produces a $\delta(t)$ term. Decompose
\begin{equation}
  G_{\mathrm{ft}}(\omega) \;=\; Q_{\mathrm{poly}}(\ii\omega)
  \;+\; G_{\mathrm{proper}}(\omega),
  \qquad
  G_{\mathrm{proper}} \xrightarrow{\;\omega\to\infty\;} 0,
\end{equation}
and the polynomial part transforms to derivatives of $\delta(t)$. For the common
case where that polynomial part is a constant matrix,
\begin{equation}
  D_\delta[i,j] \;=\; \lim_{\omega\to\infty} G_{\mathrm{ft}}[i,j](\omega).
  \label{eq:Ddelta}
\end{equation}

\begin{defn}
A rational function is \term{strictly proper} if $\deg(\text{numerator}) <
\deg(\text{denominator})$; it vanishes at infinity and contributes no $\delta(t)$.
A \term{non-proper} rational has a polynomial part and does contribute one. The
matrix $D_\delta$ collects the $\delta(t)$ coefficients of the time-domain
propagator. Physically, a nonzero entry is an \emph{immediate, same-time}
response --- in the \msrjd{} action, a response-field source $\tilde n$ at time
$t$ produces a $\delta n$ at the \emph{same} $t$.
\end{defn}

Concretely, per entry $G_{\mathrm{ft}}[i,j] = P_{ij}/Q$:
\[
  \deg(P_{ij}) < \deg(Q) \;\Rightarrow\; D_\delta[i,j] = 0,
  \qquad
  \deg(P_{ij}) = \deg(Q) \;\Rightarrow\; D_\delta[i,j] =
  \frac{\operatorname{lc}(P_{ij})}{\operatorname{lc}(Q)},
\]
the ratio of leading coefficients. The engine computes this three different ways
depending on which path produced the propagator:

\begin{itemize}
  \item In the \emph{exact path}, it is the leading-coefficient ratio over the
        exact polynomials (\file{pipeline/\_propagator.py:283}).
  \item In the \emph{lean symbolic} block, it is \code{\_omega\_inf\_limit\_fast}
        applied per entry (\file{pipeline/\_propagator.py:719}).
  \item In the \emph{rich} numerical branch, it is a \term{two-probe scaling test}
        (\file{pipeline/\_propagator.py:1150}): evaluate $G(\omega)$ at $\omega =
        10^8$ and $\omega = 10^{10}$; an entry whose value is the same at both
        probes is a constant (so it becomes a $\delta$), while one that scales like
        $1/\omega$ has decayed to zero.
\end{itemize}

\code{\_omega\_inf\_limit\_fast} (\file{pipeline/\_propagator.py:460}) deserves a
look, because it exists purely to avoid a slow tool:

\begin{lstlisting}
def _omega_inf_limit_fast(expr, omega_var):
    expr = SR(expr)
    try:
        num = expr.numerator()
        den = expr.denominator()
        num_deg = int(num.degree(omega_var))
        den_deg = int(den.degree(omega_var))
    except (AttributeError, TypeError, ValueError):
        return None
    if num_deg < den_deg:
        return SR(0)
    if num_deg > den_deg:
        return None
    # Same degree: leading-coefficient ratio.
    ...
\end{lstlisting}

\begin{gotcha}
\textbf{Sage's \code{limit} is avoided for the $\omega \to \infty$ limit.} It
routes through Maxima (a Lisp computer-algebra system bundled with Sage) and, on
raw cofactor or determinant entries, triggers thousands of FLINT divide-by-zero
warnings and Singular thrashing. \code{\_omega\_inf\_limit\_fast} does the
textbook rational-limit recipe --- compare numerator and denominator degrees,
return the leading-coefficient ratio --- entirely by hand, with no Maxima
(\file{pipeline/\_propagator.py:464}).
\end{gotcha}

\section{The time-domain propagator \texorpdfstring{$G_R(t)$}{G\_R(t)}}

Putting the pieces together --- the inverse Fourier transform of a sum of simple
poles plus a constant --- gives the time-domain \term{retarded propagator}:
\begin{equation}
  G_R[i,j](t) \;=\; D_\delta[i,j]\,\delta(t)
  \;+\; \Theta(t)\sum_k C_k[i,j]\,\ee^{\ii\omega_k t}.
  \label{eq:GR}
\end{equation}
The Heaviside step $\Theta(t)$ enforces \emph{retardation} (response only after
the cause); with $\mathrm{Im}(\omega_k) > 0$ each mode decays for $t > 0$, so the
product is well-defined.

This is exactly what \code{build\_G\_t\_matrix}
(\file{msrjd/integration/time\_domain/propagator\_td.py:135}) assembles. It
returns a dict with a \code{smooth} matrix (the pole-residue sum, \emph{without}
the $\Theta$ --- the caller applies that), a \code{delta} matrix (the $\delta(t)$
coefficients), and the time variable:

\begin{lstlisting}
smooth = matrix(SR, n_propagator, n_propagator, 0)
for k in range(len(pole_vals)):
    smooth = smooth + C_mats[k] * exp(I * pole_vals[k] * t_var)
\end{lstlisting}

Note the accumulator is initialised as a proper $n_f \times n_f$ zero matrix
\emph{before} the loop, so that even when \code{pole\_vals} is empty the result is
still a matrix (a bare \code{sum} over an empty generator would collapse to the
Python integer \code{0} and break every downstream \code{.dimensions()} call ---
the comment at \file{propagator\_td.py:252} explains the trap). When
\code{pole\_vals} \emph{is} empty, the function also issues a warning
(\file{propagator\_td.py:271}), because then only the $\delta$-part survives and
that usually means the characteristic-polynomial solve found no retarded roots.

\subsection{The complex-comparison pitfall}

Before substituting numbers, \code{build\_G\_t\_matrix} normalises the poles and
residues with a small helper, \code{\_to\_sr\_ab}
(\file{propagator\_td.py:225}):

\begin{lstlisting}
def _to_sr_ab(value):
    """Normalise only raw complex / CDF scalars; leave SR alone."""
    if isinstance(value, complex):
        return SR(float(value.real)) + SR(float(value.imag)) * I
    if isinstance(value, ComplexDoubleElement):
        c = complex(value)
        return SR(float(c.real)) + SR(float(c.imag)) * I
    # Symbolic SR expression, Sage number, etc. -- preserve exactly.
    return SR(value)
\end{lstlisting}

\begin{gotcha}
\textbf{\code{\_to\_sr\_ab} must run before \code{.subs(num\_params)}, and it must
\emph{not} cast genuine SR expressions through \code{complex()}.} The first half
of that rule: a raw Python \code{complex} fed to \code{SR(\ldots)} becomes an
opaque GiNaC node, and when that node later propagates into a product or
expansion, GiNaC's internal term-sort tries to compare two \code{complex} objects
with \code{<} and raises \code{TypeError: '<' not supported between instances of
'complex' and 'complex'}. Rewriting it as $a + b\,\ii$ with real \code{float}
components sidesteps that. The second half: an \emph{exact} closed-form pole such
as $\ii(1 - \sqrt6/10)$ must keep its algebraic content; casting it through
\code{complex()} would silently collapse it to double precision and bias
correlators that compound the error across many $\tau$ evaluations
(\file{propagator\_td.py:204}). So \code{\_to\_sr\_ab} normalises only raw complex
scalars and leaves true SR expressions untouched.
\end{gotcha}

\subsection{The delta coefficients}

\code{build\_G\_t\_matrix} prefers the precomputed \code{D\_delta}, applying
\code{num\_params} if given, and only falls back to Sage's symbolic \code{limit}
when no \code{D\_delta} was precomputed (\file{propagator\_td.py:311}):

\begin{lstlisting}
D_precomputed = propagator_data.get('D_delta')
if D_precomputed is not None:
    if num_params:
        delta_coeffs = D_precomputed.apply_map(
            lambda e: SR(e).subs(num_params) if not SR(e).is_zero() else SR(0)
        )
    else:
        delta_coeffs = D_precomputed
else:
    # Compute from G_ft via symbolic limit.
    ...
\end{lstlisting}

Since every path in phase~4 fills \code{D\_delta}, the symbolic-limit fallback is
rarely exercised.

\subsection{The per-edge lookups and the transpose}

The diagram integrators do not consume the matrix directly; they ask for one
entry at a time via \code{G\_t\_entry} (\file{propagator\_td.py:379}) for the
smooth part and \code{G\_t\_delta\_coeff} (\file{propagator\_td.py:433}) for the
$\delta$ coefficient. The smooth lookup is where the index transpose is applied:

\begin{lstlisting}
if isinstance(G_t_obj, dict):
    G_t_matrix = G_t_obj['smooth']
else:
    G_t_matrix = G_t_obj

t_var = _infer_time_variable(G_t_matrix)
entry = SR(G_t_matrix[phys_idx, resp_idx])
if t_var is not None:
    entry = entry.subs({t_var: t_expr})
if include_heaviside:
    entry = entry * heaviside(t_expr)
return entry
\end{lstlisting}

\begin{gotcha}
\textbf{The kernel layout is [response, physical], but the physical retarded
propagator is read with the indices transposed.} The physical object is ``the
response of physical field $j$ to a response-field source $i$,'' written $G^R_{j
\leftarrow i}$, which is the matrix entry $G[j, i]$ --- physical row, response
column. So \code{G\_t\_entry(phys\_idx=j, resp\_idx=i, \ldots)} reads
\code{smooth[phys\_idx, resp\_idx]} to apply the transpose. Every consumer must
use the same transpose; this one matches \code{\_get\_propagator\_entry} in
\file{msrjd/integration/symbolic.py} (\file{propagator\_td.py:44}). For an edge
$u \to v$ the time argument is $t_{\mathrm v} - t_{\mathrm u}$.
\end{gotcha}

The \code{include\_heaviside} flag multiplies in $\Theta(t)$ to enforce
retardation; only the \emph{smooth} part is returned, because the $\delta(t)$ part
is handled separately by the tree integrator (where a $\delta$ edge collapses one
integration variable). The companion \code{G\_t\_delta\_coeff}
(\file{propagator\_td.py:433}) returns that $\delta$ coefficient as a Python
\code{complex}, demoting it to a real \code{float} when the imaginary part is
negligible, and returning \code{0.0+0.0j} when there is no $\delta$ component.

\begin{gotcha}
\textbf{$\Theta(0)$ semantics differ by path.} The numerical integrator treats
$\Theta(0) = 0$ --- a zero time difference is strictly excluded from the retarded
support, enforced by a Heaviside-filter integrand wrapper in
\file{final\_integral.py} and by strict polytope inequalities. But Sage's
\emph{symbolic} \code{heaviside(0)} returns $1/2$ by default. That $1/2$ is used
only in the symbolic-display path, never in the JIT-compiled numerical integrand
(\file{propagator\_td.py:30}). The two coexist precisely because they never meet.
\end{gotcha}

\subsection{Inferring the time variable without triggering Maxima}

A small but instructive helper, \code{\_infer\_time\_variable}
(\file{propagator\_td.py:459}), recovers the symbolic $t$ from the smooth matrix.
Once parameters are substituted, $t$ is the only free variable left in each
entry. The naive way to find a nonzero entry would be \code{entry == 0} --- but
that comparison routes to Maxima, which fails on entries carrying embedded Python
\code{complex} coefficients. So the helper checks \code{entry.variables()}
directly and only uses \code{is\_trivial\_zero()} for scalar entries. It is a tiny
function, but it embodies a recurring theme in this subsystem: \emph{avoid the
slow, fragile symbolic tool when a cheap structural check will do.}

\section{A concrete example: the OU propagator}

It is worth grounding all of this in the simplest model, the white-noise OU
process whose free action is $S_{\mathrm{free}} \supset \tilde n(\Dt n + \mu n)$.
There is one field, so $n_f = 1$ and every matrix is $1\times1$:

\begin{itemize}
  \item The single kernel entry is $K[\tilde n, n] = \Dt + \mu$, which
        \code{\_to\_kernel} rewrites as $\delta' + \mu\,\delta$.
  \item Under the Fourier convention \eqref{eq:ftconv}, $K_{\mathrm{ft}} = \ii\omega
        + \mu$.
  \item The propagator is $G_{\mathrm{ft}} = 1/(\ii\omega + \mu)$, with a single
        pole at $\omega = \ii\mu$ --- which has $\mathrm{Im} > 0$ for $\mu > 0$, so
        it is retarded.
  \item The residue of $G_{\mathrm{ft}}$ at the pole is $1/\ii = -\ii$, so by the
        $C_k$ convention \eqref{eq:Ck} the residue matrix is $C_0 = \ii\cdot(-\ii)
        = 1$.
  \item $G_{\mathrm{ft}}$ is strictly proper ($\deg P = 0 < 1 = \deg Q$), so
        $D_\delta = 0$.
\end{itemize}

Assembling \eqref{eq:GR}: $G_R(t) = \Theta(t)\cdot 1\cdot\ee^{\ii(\ii\mu)t} =
\Theta(t)\,\ee^{-\mu t}$ --- the textbook retarded relaxation. A two-field system
(a density coupled to a potential, say) produces a $2\times2$ $K_{\mathrm{ft}}$
with two poles, and a $\tilde n \leftrightarrow \delta n$ coupling that yields a
nonzero $D_\delta$ entry --- an instantaneous $\delta(t)$ piece in $G_R$.

\section{Robustness machinery and external tools}

This subsystem leans almost entirely on \term{SageMath} and, through it, several
computer-algebra back-ends. None of nauty, networkx, numba, or sympy is called
\emph{directly} here --- though, as noted, Sage's \code{fourier\_transform} routes
through SymPy internally. A consolidated tour of the tools, since this is where
they all first appear:

\begin{description}
  \item[SageMath / \code{SR}] The Symbolic Ring (GiNaC/Pynac-backed) and the glue
        that routes each operation to the right back-end. \code{matrix(ring, data)}
        builds matrices over \code{SR}, \code{CDF}, or an exact field; matrix
        methods used include \code{.inverse()}, \code{.det()}, \code{.apply\_map(fn)}
        (entrywise function application), and \code{.dimensions()}.
  \item[\code{QQ}, \code{CDF}, \code{CyclotomicField(4)}] The rationals
        $\mathbb{Q}$, the machine-precision complex field, and $\mathbb{Q}(i)$ ---
        the three number fields the residue paths move between.
  \item[\code{PolynomialRing} / fraction field] $R[\omega]$ and
        $\mathrm{Frac}(R[\omega])$; the latter reduces each inverse entry to
        canonical irreducible $P/Q$ form when $R$ is an exact field.
  \item[NumPy] Imported lazily inside the numerical branches; \code{np.linalg.inv}
        / \code{np.linalg.det} for the matrix evaluations, \code{np.delete} to form
        the minors of the cofactor adjugate, and \code{np.polyfit} for the
        fallback pole-finder. NumPy is preferred for residues because its LU-based
        linear algebra avoids the sum-of-rationals cancellation the symbolic path
        suffers near poles.
  \item[\code{cysignals}] A small Cython library (bundled with Sage) that can
        interrupt or recover from signals raised inside native C/Cython routines.
        Plain Python \code{try/except} cannot catch a hang or crash deep inside
        Singular; \code{cysignals} can.
  \item[Python \code{signal}] The POSIX signal API. \code{\_run\_with\_timeout}
        uses \code{SIGALRM} + \code{signal.alarm} to wall-clock-bound the exact
        inverse and the symbolic-inverse block, degrading to a plain call where
        \code{SIGALRM} is absent.
\end{description}

The most colourful piece of robustness machinery is the \code{factor()} hazard.
\code{build\_propagator} deliberately returns the \emph{raw}, unfactored inverse,
because \code{factor()} can hang or segfault Singular (Sage's polynomial CAS). A
separate, purely cosmetic helper \code{factor\_propagator}
(\file{pipeline/\_propagator.py:366}) applies factoring only for report display,
and it does so with \emph{two} layers of protection (\code{\_safe\_factor},
\file{pipeline/\_propagator.py:422}):

\begin{lstlisting}
from cysignals.alarm import alarm, cancel_alarm, AlarmInterrupt
try:
    alarm(timeout)
    try:
        return e.factor()
    finally:
        cancel_alarm()
except (RuntimeError, ValueError, TypeError, ArithmeticError,
        AlarmInterrupt):
    return e
except BaseException as exc:
    if exc.__class__.__name__ == 'SignalError':
        return e
    raise
\end{lstlisting}

The \code{alarm(timeout)} arms a timer that raises \code{AlarmInterrupt} to stop a
runaway \code{factor()} loop, and on any failure the helper falls back to the
unfactored entry. \code{factor\_propagator} adds a one-shot ``bail'' flag that
permanently skips factoring after any single entry exceeds a threshold, because
the alarm does not always interrupt Singular's native loop.

\begin{gotcha}
\textbf{The \code{cysignals} \code{SignalError} is caught by class name, not by
\code{isinstance}.} In some Sage builds \code{SignalError} (Sage's wrapper around
a native segfault) is \emph{not} a subclass of \code{Exception}. So both
\code{\_safe\_factor} and \code{build\_propagator} catch \code{BaseException} and
check \code{exc.\_\_class\_\_.\_\_name\_\_ == 'SignalError'}, re-raising everything
else --- which is what keeps \code{KeyboardInterrupt} and \code{SystemExit} from
being swallowed (\file{pipeline/\_propagator.py:451}, \file{:746}).
\end{gotcha}

\section{Caching}

The symbolic stages of the propagator are parameter-independent, so they are
cached to disk. \code{build\_propagator} slugifies the model name into a tag and
opens a \code{PipelineCache} on that directory
(\file{pipeline/\_propagator.py:547}):

\begin{lstlisting}
prop_tag = re.sub(r'[^A-Za-z0-9]+', '_', model['name']).strip('_').lower()
cache_dir = f"{cache_dir_root}/{prop_tag}"
cache = PipelineCache(cache_dir)
\end{lstlisting}

\code{PipelineCache} (\file{msrjd/core/cache.py}) wraps Sage's pickling ---
\code{save} writes a \code{.sobj} (Sage object) file, \code{load} reads it back ---
so the cache lands at \file{saved\_theories/<tag>/propagator.sobj}. The cache has
two staleness guards (\file{pipeline/\_propagator.py:551}): it rebuilds if the
cached \code{nf} differs from \code{ft.\_n\_tilde} (catching a field-list edit that
did not change the model name), and it rebuilds if the model is spatial but the
cached propagator predates the spatial block (no \code{G\_tx\_sym}).

\begin{gotcha}
\textbf{Closures do not pickle, so the spatial \code{G\_tx} callables are attached
\emph{after} the cache save.} For a spatial model, \code{build\_propagator} builds
the heat-kernel propagator block and merges it into \code{prop}, but the runtime
callables \code{G\_tx} are functions, which Sage cannot pickle. So the cache is
saved \emph{first} (\file{pipeline/\_propagator.py:817}), with only the picklable
symbolic data, and the \code{G\_tx} callables are rebuilt and attached afterwards
via \code{make\_g\_tx\_callables} (\file{pipeline/\_propagator.py:830}). On a cache
hit the same rebuild step runs. The spatial heat-kernel propagator is the subject
of its own chapter; here it appears only at its caching seam.
\end{gotcha}

\section{Known limitations and failure modes}

A short field guide to the ways this subsystem can refuse or mislead, gathered in
one place:

\begin{description}
  \item[Higher-multiplicity (Jordan-block) poles are not handled.] A
        multiplicity-$\ge 2$ pole produces a $\tau^k\,\ee^{-pt}$ mode that the
        single-pole residue formula \eqref{eq:residue} cannot represent. The exact
        path \emph{bails} on a non-squarefree denominator
        (\file{pipeline/\_propagator.py:196}); the numerical path \emph{silently
        leaves a zero residue} where $|Q'| < 10^{-30}$. A future version would need
        the $m$-th-derivative residue formula \emph{and} a downstream integrator
        that can absorb the polynomial-times-exponential modes. This is the most
        significant known limitation.
  \item[Non-rationalisable parameters knock out the exact path.] An irrational
        symbolic parameter throws inside \code{\_sr\_complex\_to\_CF}; the exact
        tier returns \code{None} and the numerical cofactor path takes over.
  \item[Surviving free symbols give zero poles.] If a kernel symbol is
        unsubstituted or a saddle value is missing from \code{num\_params}, the
        characteristic polynomial has non-numeric coefficients and the root-finder
        returns nothing --- surfaced as a warning, then an empty pole list.
  \item[An empty pole list makes the smooth propagator identically zero.] Only
        $D_\delta$ then contributes. \code{build\_G\_t\_matrix} warns about this
        (\file{propagator\_td.py:271}) precisely so an unexpected
        no-retarded-roots solve does not slip through silently.
\end{description}

\section{What you just saw}

The propagator is two phases of the trace --- \texttt{[2/7]} and \texttt{[4/7]}
--- and one object, $\Gret$, built in three movements:

\begin{enumerate}
  \item \textbf{Action to frequency domain.} Read the kernel matrix $K$ out of the
        bilinear free action, rewrite each entry in $\delta/\delta'$ kernel form,
        Fourier transform under $\ee^{-\ii\omega t}$, and invert to get
        $G_{\mathrm{ft}}(\omega)$ --- with a matrix-size gate that skips symbolic
        inversion for large matrices and defers them to the numerical stage. All of
        this is parameter-independent and cached.
  \item \textbf{Poles and residues.} Substitute the numeric parameters; find the
        retarded poles ($\mathrm{Im}(\omega) > 0$) and the residue matrices $C_k =
        \ii\,\mathrm{Res}\,G_{\mathrm{ft}}$, preferring an exact $\mathbb{Q}(i)$
        path and falling back to a numerical cofactor path, with residues always
        evaluated numerically to dodge catastrophic cancellation.
  \item \textbf{Back to the time domain.} Assemble $G_R(t) = D_\delta\,\delta(t) +
        \Theta(t)\sum_k C_k\,\ee^{\ii\omega_k t}$, and serve it one edge at a time
        --- transposed to [physical, response] --- to the diagram integrators.
\end{enumerate}

With $\Gret$ in hand, the engine has the building block for every Feynman edge.
The next chapters open up what is wired together with those edges: the mean-field
saddle the parameters are evaluated at --- which is where \code{num\_params} comes
from --- the diagram enumeration that decides \emph{which} graphs to build, and
the loop integration that finally turns a graph of $\Gret$ edges into a number.
